\lyseis{7.10}{Τυποποιημένη μορφή μεγάλων και μικρών αριθμών}
\sol{1} α) $5\cdot10^{-4}$ \quad β) $7{,}2\cdot10^{-5}$ \quad γ) $3\cdot10^{-1}$ \quad δ) $4{,}5\cdot10^{4}$ \quad ε) $8\cdot10^{-7}$ \quad στ) $6{,}3\cdot10^{6}$.
\sol{2} α) $0{,}000032$ \quad β) $0{,}0091$ \quad γ) $0{,}4$ \quad δ) $2\,500\,000$ \quad ε) $0{,}000101$ \quad στ) $7$.
\sol{3} α) $3\cdot10^{-3}$ \quad β) $1{,}2\cdot10^{-3}$ \quad γ) είναι ήδη τυποποιημένο \quad δ) $1\cdot10^{3}$.
\sol{4} α) $>$ ($-5>-6$) \quad β) $<$ ($-3<-2$) \quad γ) $>$ ($4>-4$) \quad δ) $>$ ($3{,}1>3{,}09$ με ίδιο εκθέτη).
\sol{5} $5\cdot10^{-5}<1{,}2\cdot10^{-3}<3\cdot10^{-3}<9\cdot10^{-1}<2\cdot10^{2}$.
\sol{6} α) $6\cdot10^{2}$ \quad β) $3\cdot10^{10}$ \quad γ) $20\cdot10^{2}=2\cdot10^{3}$ \quad δ) $2\cdot10^{-3}$.
\sol{7} $0{,}00045$, $25\,000\,000$ και $0{,}000\,000\,007$.
\sol{8} α) Σωστό. β) Λάθος: $1{,}2\cdot10^{-3}$. γ) Λάθος: το $10^{-5}$ είναι μεγαλύτερο. δ) Σωστό. ε) Σωστό. στ) Λάθος: $5{,}5\cdot10^{-3}<5{,}5\cdot10^{-2}$.
\sol{9} $10^{-7}<2\cdot10^{-6}<7\cdot10^{-6}<7\cdot10^{-5}$ (ιός, βακτήριο, αιμοσφαίριο, τρίχα). Η τρίχα είναι $10$ φορές χοντρότερη από το αιμοσφαίριο.
\sol{10} $9{,}1\cdot10^{-31}$ έχει $30$ μηδενικά μετά την υποδιαστολή πριν από το $9$ ($0{,}000\ldots09\ldots$ με $30$ μηδενικά).
\sol{11} $\frac{1{,}5\cdot10^{11}}{3\cdot10^{8}}=5\cdot10^{2}=500$ s, δηλαδή περίπου $8{,}3$ λεπτά.
\sol{12} $2\cdot10^{-9}$ μ.
\sol{13} $500\cdot10^{-4}=5\cdot10^{-2}$ μ. $=5$ εκατοστά.
\sol{14} $10^{6}$, δηλαδή ένα εκατομμύριο κόκκοι.
\sol{15} Γιατί για να πάμε στον αριθμό από το $\alpha\ge1$ διαιρούμε με δύναμη του $10$ (μετακινούμε την υποδιαστολή αριστερά), και αυτό είναι ο αρνητικός εκθέτης $10^{-\nu}=\frac{1}{10^\nu}$.
\sol{16} $1=1\cdot10^{0}$ και $10=1\cdot10^{1}$.
\sol{17} Δεκαδικά: $0{,}002>0{,}0005$. Εκθέτες: $-3>-4$. Άρα $2\cdot10^{-3}>5\cdot10^{-4}$.
\sol{18} $3\cdot10^{-4}\cdot10^{6}=3\cdot10^{2}=300$.
\sol{19} Γιατί ο $12$ δεν είναι μικρότερος από το $10$. Γράφεται $1{,}2\cdot10^{-3}$.
\sol{20} Ελεύθερη απάντηση. Π.χ. άτομο $\approx10^{-10}$ μ., ιός $\approx10^{-7}$ μ., βακτήριο $\approx2\cdot10^{-6}$ μ.
